\section{Powertrain Modeling and Motor Selection}

\begin{multicols}{2}
The foundation of the electric vehicle (EV) powertrain begins with accurately quantifying the opposing environmental forces to derive the requisite mechanical output of the motor. This section details the fundamental longitudinal vehicle dynamics, motor rating selection, and the implementation of advanced control algorithms required to manage the tractive power effectively.

\subsection{Longitudinal Vehicle Dynamics}

To define the continuous and peak operational limits of the electric motor, the total tractive effort ($F_{te}$) required at the tire contact patch must be explicitly calculated. This effort must overcome four distinct environmental and inertial components: Aerodynamic Drag ($F_{ad}$), Rolling Resistance ($F_{rr}$), Gradient Force ($F_{g}$), and Acceleration Force ($F_{a}$).

Assuming a standard lightweight EV geometry, the vehicle mass ($m$) is established at 300 kg, possessing a frontal area ($A$) of 1.0 m$^2$, and a target top speed ($v$) of 80 km/h (22.22 m/s). 

\textbf{1. Aerodynamic Drag ($F_{ad}$):} As velocity increases, the required force to displace the fluid volume scales quadratically. Given an air density ($\rho$) of 1.225 kg/m$^3$ and a conservative drag coefficient ($C_d$) of 0.4:
\begin{equation} \label{eq:aero_drag}
\begin{split}
F_{ad} &= \frac{1}{2} \rho v^2 C_d A \\
&= 0.5 \cdot 1.225 \cdot (22.22)^2 \cdot 0.4 \cdot 1.0 \\
&\approx 121.0\,\mathrm{N}
\end{split}
\end{equation}

\textbf{2. Rolling Resistance ($F_{rr}$):} The hysteresis losses generated by the deformation of the viscoelastic tire compound against the tarmac are quantified by the coefficient of rolling resistance ($\mu_r \approx 0.015$). With gravitational acceleration ($g$) at 9.81 m/s$^2$:
\begin{equation} \label{eq:rolling_res}
\begin{split}
F_{rr} &= m \cdot g \cdot \mu_r \\
&= 300 \cdot 9.81 \cdot 0.015 \\
&= 44.1\,\mathrm{N}
\end{split}
\end{equation}

\textbf{3. Gradient/Hill Climbing Force ($F_{g}$):} To ensure operational viability on inclines, the powertrain must overcome the gravitational component acting parallel to the road surface. Assuming a design operational gradient ($\alpha$) of $5^\circ$ ($\sim$8.7\% grade):
\begin{equation} \label{eq:gradient_force}
\begin{split}
F_{g} &= m \cdot g \cdot \sin(\alpha) \\
&= 300 \cdot 9.81 \cdot \sin(5^\circ) \\
&= 256.3\,\mathrm{N}
\end{split}
\end{equation}

\textbf{4. Acceleration Force ($F_{a}$):} To achieve a target 0-80 km/h sprint in 11 seconds, an average linear acceleration ($a$) of approximately 2 m/s$^2$ is required:
\begin{equation} \label{eq:accel_force}
\begin{split}
F_{a} &= m \cdot a \\
&= 300 \cdot 2.0 \\
&= 600.0\,\mathrm{N}
\end{split}
\end{equation}

The cumulative tractive effort ($F_{te\_peak}$) under maximum acceleration on a $5^\circ$ incline is the direct summation of these vectors:
\begin{equation} \label{eq:total_tractive}
\begin{split}
F_{te\_peak} &= F_{ad} + F_{rr} + F_{g} + F_{a} \\
&= 121.0 + 44.1 + 256.3 + 600.0 \\
&= 1021.4\,\mathrm{N}
\end{split}
\end{equation}
\end{multicols}

\begin{figure}[H]
\centering
\includegraphics[width=\textwidth]{"Figures Images/figure6_foc.png"}
\caption{Field Oriented Control (FOC) Block Diagram utilized for precise regulation of IPM instantaneous torque.}
\label{fig:foc_diagram}
\end{figure}

\begin{multicols}{2}
\subsection{Motor Rating Determination}

Translating the linearly derived tractive effort into rotational motor requirements demands accounting for the mechanical drivetrain efficiency ($\eta \approx 0.90$) and the dynamic tire radius ($r = 0.25$ m).

The absolute peak power ($P_{peak\_motor}$) required at the motor shaft during maximum velocity acceleration is:
\begin{equation} \label{eq:motor_peak_pwr}
\begin{split}
P_{peak\_motor} &= \frac{F_{te\_peak} \cdot v}{\eta} \\
&= \frac{1021.4 \cdot 22.22}{0.90} \\
&\approx 25,216\,\mathrm{W} \approx 25.2\,\mathrm{kW}
\end{split}
\end{equation}

Conversely, the nominal continuous power rating ($P_{cont\_motor}$) is defined by the steady-state cruise condition on flat ground ($F_g = 0$, $F_a = 0$):
\begin{equation} \label{eq:motor_cont_pwr}
\begin{split}
P_{cont\_motor} &= \frac{(F_{ad} + F_{rr}) \cdot v}{\eta} \\
&= \frac{(121.0 + 44.1) \cdot 22.22}{0.90} \\
&\approx 4,076\,\mathrm{W} \approx 4.1\,\mathrm{kW}
\end{split}
\end{equation}

Therefore, the finalized motor specification dictates a continuous rating of 5 kW to maintain a thermal safety factor, with a peak transient overload capability of 25.2 kW.

\subsection{Motor Selection: IPM vs. SPM}

The selection between a Surface Permanent Magnet (SPM) and Interior Permanent Magnet (IPM) motor dramatically influences the torque density and high-speed field-weakening capabilities of the 72V powertrain.

While an SPM motor simplifies the rotor construction by adhering the magnets to the outer rotor circumference, it suffers from a near-homogenous magnetic reluctance path ($L_d \approx L_q$), restricting it exclusively to magnetic alignment torque. Furthermore, the fragile mechanical retention limits rotational speeds, reducing power density.

Therefore, an \textbf{Interior Permanent Magnet (IPM)} topology is explicitly selected for this powertrain. Embedding the magnets within the rotor lamination steel creates magnetic saliency ($L_q > L_d$). This asymmetric reluctance path unlocks a secondary, profound torque-producing component. The total electromagnetic torque ($T_e$) of the selected IPM motor is defined as:

\begin{equation} \label{eq:ipm_torque}
T_e = \underbrace{\frac{3}{2} p \lambda_f i_q}_{\text{Magnetic Torque}} + \underbrace{\frac{3}{2} p (L_d - L_q) i_d i_q}_{\text{Reluctance Torque}}
\end{equation}

where $p$ represents pole pairs, $\lambda_f$ is the permanent magnet flux linkage, and $i_d, i_q$ are the direct and quadrature axis stator currents. At lower base speeds, the controller biases current heavily towards $i_q$ to maximize magnetic alignment. As velocity increases, the controller manipulates $i_d$ (the negative $d$-axis current), strategically leveraging the $(L_d - L_q)$ reluctance inequality to generate supplementary torque, enabling superior overall system efficiency and high-speed performance compared directly against an SPM equivalent.

\subsection{Matching Control Method (FOC)}

Unlocking the complex reluctance torque of the IPM architecture necessitates an advanced Field Oriented Control (FOC) strategy, fundamentally superior to primitive Six-Step BLDC commutation or scalar V/Hz methods.

The FOC algorithm mathematically decouples the 3-phase AC stator currents ($i_a, i_b, i_c$) via consecutive Clarke and Park transformations into a synchronous rotating reference frame aligned with the internal rotor magnet flux. This reduces the complex AC equations into two controllable DC orthogonal components: flux-producing current ($i_d$) and torque-producing current ($i_q$).

\textbf{Maximum Torque Per Ampere (MTPA) Strategy:}
Unlike an SPM motor where MTPA strictly enforces $i_d = 0$, optimal control of the selected IPM motor requires a dynamic $i_d$ reference to maximize the combined magnetic and reluctance torque output per unit of stator current. The MTPA algorithm calculates the optimal negative $i_d$ trajectory as a function of the operational torque demand:

\begin{equation} \label{eq:ipm_mtpa}
i_d^* = \frac{\lambda_f}{2(L_q - L_d)} - \sqrt{ \frac{\lambda_f^2}{4(L_q - L_d)^2} + (i_q^*)^2 }
\end{equation}

Translation of the resultant Proportional-Integral (PI) controller reference voltages back into physical gating commands is achieved via Space Vector PWM (SVPWM). SVPWM synthesizes the rotating voltage vector by symmetrically modulating the eight fundamental switching states of the 3-phase inverter bridge. This matching topology provides a 15\% superior utilization of the 72V DC bus compared to standard sinusoidal PWM, critical for extracting the full 25.2 kW peak rating from a low-voltage battery pack.
\end{multicols}
